%==============================================================================
%  ANNEXES
%==============================================================================
% Matériel opérationnel conservé hors du corps : tuteur conversationnel,
% instruments d'enquête.
%==============================================================================

\chapter{Tuteur conversationnel}
\label{annexe:templates}

Les vingt patrons A1--D5, y compris leurs consignes, figurent au
tableau~\ref{tab:templates} du chapitre~\ref{chap:conception}. Cette
annexe ne conserve que le réglage du tuteur, distinct de la cascade
de génération.

Le tuteur reprend le même identifiant de patron que la génération, mais
avec un prompt système distinct, une température \emph{par état} et un
historique borné aux dix derniers tours. Le contenu adaptatif déjà
affiché est injecté comme premier message, afin d'éviter la recopie
mot à mot.

\paragraph{Températures par état.}
Le tableau~\ref{tab:annexe-temp-chat} présente le protocole du tuteur conversationnel.

{\small
\renewcommand{\arraystretch}{1.18}
\setlength{\tabcolsep}{4pt}
\begin{xltabular}{\textwidth}{|>{\raggedright\arraybackslash}X|c|l|}
\hline
\rowcolor{gray!20}
\textbf{État} & \textbf{$T$} & \textbf{Consigne émotionnelle (extrait)} \\
\hline
\endfirsthead
\hline
\rowcolor{gray!20}
\textbf{État} & \textbf{$T$} & \textbf{Consigne émotionnelle (extrait)} \\
\hline
\endhead
\endfoot
\endlastfoot
Confusion & 0{,}25 & Décomposer en 2--3 micro-étapes, question de vérification. \\
\hline
Frustration & 0{,}35 & Valider l'effort, une seule action de déblocage. \\
\hline
Engagement & 0{,}45 & Approfondir, mini-défi, pensée critique. \\
\hline
Ennui & 0{,}55 & Cas réel, anecdote, éviter les longues définitions. \\
\hline
\end{xltabular}
\xlcaption{Réglages du tuteur conversationnel.}
\label{tab:annexe-temp-chat}
}


\paragraph{Règles absolues du message système (extrait).}
Réponses uniquement en français, 3 à 8 phrases par défaut (jusqu'à 12 si
l'apprenant demande un exemple ou un exercice), s'appuyer sur le RAG
lorsqu'il est fourni, ne pas inventer de faits techniques, rester
focalisé sur le concept en cours.

%==============================================================================
\chapter{Questionnaires et protocoles d'évaluation}
\label{annexe:protocoles}







\section{Enquête d'idée (exécutée, $n=15$)}
\label{sec:annexe-enquete-idee}

Quinze étudiants ont répondu
(\texttt{https://questionnaire-etudiants.netlify.app}, 
source : \texttt{questionnaires/etudiants.html})
sans ouvrir le prototype. La page décrit EmoLearn en trois
couches, puis pose les items au conditionnel. Résultats :
tableau~\ref{tab:enquete-idee-n15}, 
fichier \texttt{results/forms/l4\_enquete\_idee\_reponses.csv}.

\paragraph{Consignes (extrait).}
«~Vous n'utilisez pas le prototype. Répondez seulement d'après la
description, comme si cette idée existait dans un cours en ligne.~»
Consentement obligatoire, la page ne filme pas et n'enregistre ni
webcam ni micro.

Le tableau~\ref{tab:annexe-enquete-idee} présente les items de l'enquête d'idée.

{\small
\renewcommand{\arraystretch}{1.18}
\setlength{\tabcolsep}{4pt}
\begin{xltabular}{\textwidth}{|>{\bfseries}l|>{\raggedright\arraybackslash}X|c|}
\hline
\rowcolor{gray!20}
\textbf{Code} & \textbf{Item (libellé du formulaire)} & \textbf{Éch.} \\
\hline
\endfirsthead
\hline
\rowcolor{gray!20}
\textbf{Code} & \textbf{Item (libellé du formulaire)} & \textbf{Éch.} \\
\hline
\endhead
\endfoot
\endlastfoot
Profil & Niveau d'études (texte libre). & --- \\
\hline
Familiarité & Familiarité avec les cours en ligne / MOOC
(ancres : pas du tout --- beaucoup). & 1--5 \\
\hline
Utilité & D'après la description (sans avoir testé le logiciel) :
un contenu qui s'adapte à l'engagement me semblerait utile. & 1--7 \\
\hline
Rester dans le cours & D'après la description : un tel contenu
m'aiderait à rester dans le cours. & 1--7 \\
\hline
Satisfaction & D'après la description : je serais satisfait(e) de
recevoir ce type de texte pendant un module. & 1--7 \\
\hline
Acceptation captation & D'après la description : j'accepterais
qu'une webcam / un micro estime mon engagement pendant un cours
(données non stockées à des fins commerciales). & 1--7 \\
\hline
Ouvertes & Qu'est-ce qui aiderait~? Qu'est-ce qui gênerait~?
Une suggestion. & texte \\
\hline
\end{xltabular}
\xlcaption{Enquête d'idée : profil, Likert et questions ouvertes.}
\label{tab:annexe-enquete-idee}
}

Ancrage des Likert 1--7 : 1 = pas d'accord, 7 = tout à fait d'accord.

\paragraph{SUS imaginé.} Dix items \textcite{brooke1996sus}, Likert
1--5, reformulés au conditionnel (titre du bloc : «~Facilité d'usage
imaginée~»). Consignes : «~Ces 10 questions ne portent pas sur un essai
réel. Imaginez seulement le système décrit plus haut.~» Score :
items impairs $(v-1)$, items pairs $(5-v)$, somme $\times 2{,}5$
(intervalle $0$--$100$). Ce n'est pas un SUS après usage.

Le tableau~\ref{tab:annexe-enquete-sus} présente les items du questionnaire SUS.

{\small
\renewcommand{\arraystretch}{1.18}
\setlength{\tabcolsep}{4pt}
\begin{xltabular}{\textwidth}{|>{\bfseries}l|>{\raggedright\arraybackslash}X|}
\hline
\rowcolor{gray!20}
\textbf{Item} & \textbf{Libellé du formulaire} \\
\hline
\endfirsthead
\hline
\rowcolor{gray!20}
\textbf{Item} & \textbf{Libellé du formulaire} \\
\hline
\endhead
\endfoot
\endlastfoot
SUS~1 & Je voudrais utiliser ce type de système fréquemment. \\
\hline
SUS~2 & Je le trouverais inutilement complexe. \\
\hline
SUS~3 & Je le trouverais facile à utiliser. \\
\hline
SUS~4 & J'aurais besoin d'un support technique pour l'utiliser. \\
\hline
SUS~5 & Les fonctions me sembleraient bien intégrées. \\
\hline
SUS~6 & Il y aurait trop d'incohérences. \\
\hline
SUS~7 & La plupart des gens l'apprendraient très rapidement. \\
\hline
SUS~8 & Je le trouverais lourd à utiliser. \\
\hline
SUS~9 & Je me sentirais en confiance en l'utilisant. \\
\hline
SUS~10 & Il me faudrait apprendre beaucoup de choses avant de m'en servir. \\
\hline
\end{xltabular}
\xlcaption{Enquête d'idée : SUS imaginé (10 items, 1--5).}
\label{tab:annexe-enquete-sus}
}

%==============================================================================
